\section{Limitations}
\label{sec:limitations}

The present evidence is diagnostic rather than a complete confirmatory
evaluation. First, the operating threshold was selected on the same
development traces used for reporting; a disjoint calibration/evaluation split
is needed before drawing general conclusions about the selected point. Second,
all results use one backbone, greedy decoding, and one run on the first
\(1{,}000\) development examples. The \musique{} slice contains only two-hop
questions and understates the difficulty of the complete development set.
Third, external baselines are paper-reported rather than reproduced in the same
harness. Fourth, original and decomposition-path retrieval differ in title
matching, and original-question evidence union has not been isolated in a
separate ablation. The retrieval-recall gain therefore belongs to the complete
query policy, not uniquely to decomposition.

The threshold sweep is a fixed-chain replay rather than an independent live
rerun at every point. It preserves a controlled comparison but cannot capture
all downstream changes caused by an earlier retrieval decision. Finally,
\tailleak{} measures the shape of a model distribution, not factual
correctness. It cannot flag a false answer generated with extreme confidence,
and routing cannot repair questions for which both candidates are wrong.

\section{Conclusion}
\label{sec:conclusion}

We introduced \tailleak{}, a training-free confidence statistic that isolates
probability mass leaving a model's leading token set and accumulates it across
an answer. In the evaluated Llama-3-8B traces, it ranks direct-answer errors
better than representative entropy, likelihood, disagreement, margin, and
self-evaluation signals, and it serves as a useful sub-question retrieval gate
inside \method{}. The results also set a clear boundary: most decomposed steps
still benefit from retrieval, late routing adds only a small correction, and
confident hallucinations remain. Broader backbones, held-out calibration, and
controlled retrieval ablations are the next tests required to establish how
far the signal generalizes.
\label{sec:main-end}
